\section{Expanding CNN Beyond Classification}
In the previous chapter, we have seen how CNNs can be used for image classification tasks.

CNNs can however be applied to a wide range of other tasks beyond classification, such as \textbf{object detection}, 
\textbf{semantic segmentation}, and \textbf{captioning}, among others.
In this chapter, we will explore some of these applications and how CNNs can be adapted to solve them.

\subsection{Object Detection}
\textbf{Object detection} is the task of identifying and localizing objects within an image. It can be seen as a combination of
\textbf{regression} (where is the object located)and \textbf{classification} (what object is present).
The task input would be an image, and the output would be a set of bounding boxes around the detected objects, along with their corresponding class labels.

This task is more complex than basic image classification, as it requires the model not only to recognize the presence of objects, 
but also to accurately localize them within the image.
This makes also the training process more challenging, as several issues can arise, such as:
\begin{itemize}
    \item \textbf{Data scarcity}: Object detection datasets are often smaller than classification datasets. Also,
    being the task more complex, it requires more data to train a model that can generalize well.
    \item \textbf{Poor localization}: the deep CNN proposed for classification may not be able to localize the objects accurately, 
    because of their large receptive field and the pooling layers that reduce the spatial resolution of the feature maps.
\end{itemize}

To address these issues, several architectures have been proposed, such as \textbf{R-CNN}, \textbf{Fast R-CNN}, and \textbf{YOLO}, among others.

\subsubsection{Before CNNs: DPM}
Before the rise of deep learning and CNNs, one of the most popular approaches for object detection was the \textbf{Deformable Part Model} (DPM).
DPM is a method that models objects as a collection of parts, where each part can be described by a set of features. 
The model learns to detect these parts and their spatial relationships to identify the presence of an object in an image.

This approach was quite effective for its time, but reached a plateau in performance already in 2010, before the advent of deep learning.

\subsubsection{Selective Search}
One of the key challenges in object detection is to efficiently search for objects in an image.
\textbf{Selective Search} is a method that generates a set of candidate regions (also called \textbf{proposals}) in an image that are likely to contain objects. 

Given the proposals, the process of object detection can be simplified to a classification task, where the model only needs to classify the proposals as containing an object or not,
which is much easier than searching for objects in the entire image.

To detect objects at any scale, Selective Search uses a hierarchical grouping of superpixels, which are small regions of an image that 
have similar features (such as color, texture, or intensity).

The algorithm starts by generating all the possible candidate objects. 
Similar regions are then recursively merged together based on their similarity. 
Once the merging process is complete, the algorithm uses the resulting regions as candidate object proposals for the object detection task.

\subsubsection{R-CNN}
\textbf{R-CNN} (Region-based Convolutional Neural Networks) is one of the first architectures that successfully applied CNNs to object detection. 
The R-CNN architecture consists of three main steps:
\begin{itemize}
    \item \textbf{Region Proposal}: The first step is to generate a set of candidate regions (proposals) in the image using Selective Search.
    
    \item \textbf{Feature Extraction}: Each proposal is then first resized to a fixed size (e.g., 224x224) and then passed through a CNN to extract a fixed-length feature vector.
    
    \item \textbf{Classification}: the extracted features are then used to classify each proposal as containing an object or not, and to assign a class label to the detected objects.
    This was typically done using a Support Vector Machine (SVM) classifier, which was trained on the features extracted from the proposals.
\end{itemize}

\begin{figure}[H]
    \centering
    \resizebox{\textwidth}{!}{\usetikzlibrary{arrows.meta, calc, positioning}

\begin{tikzpicture}[
    >=Stealth,
    flow/.style={-{Stealth[length=7pt]}, line width=1pt},
    lbl/.style={font=\small, align=center},
    steplbl/.style={font=\small\bfseries, align=center, text=dlblue!35!black},
    box/.style={draw, thick, rounded corners, align=center, font=\bfseries\small},
    propedge/.style={draw, line width=1pt, rounded corners=1pt},
]
    % ---- helper: a schematic "photo" (landscape icon) clipped to a square ----
    \def\photo#1#2#3{% #1=center coordinate name, #2=size, #3=fill
        \begin{scope}
            \clip (#1) ++(-#2/2,-#2/2) rectangle ++(#2,#2);
            \fill[#3] (#1) ++(-#2/2,-#2/2) rectangle ++(#2,#2);
            \fill[dlyellow!85!white] ($(#1)+(-0.30*#2,0.30*#2)$) circle (0.11*#2);
            \fill[dlgreen!45!black] ($(#1)+(-0.55*#2,-0.5*#2)$) -- ++(0.55*#2,0.42*#2) -- ++(0.55*#2,-0.42*#2) -- cycle;
            \fill[dlgreen!25!black] ($(#1)+(-0.05*#2,-0.5*#2)$) -- ++(0.38*#2,0.3*#2) -- ++(0.38*#2,-0.3*#2) -- cycle;
        \end{scope}
        \draw[draw=dlgray!80, thick] (#1) ++(-#2/2,-#2/2) rectangle ++(#2,#2);
    }

    \def\S{2.1}   % photo size

    % ===== Stage 1: input image =====
    \coordinate (img1) at (0,0);
    \photo{img1}{\S}{dlsky!12}
    \node[steplbl] at ($(img1)+(0,\S/2+0.45)$) {Input Image};

    % ===== Stage 2: selective search -> region proposals =====
    \coordinate (img2) at (4.6,0);
    \photo{img2}{\S}{dlsky!12}
    \node[steplbl, align=center] at ($(img2)+(0,\S/2+0.45)$) {Selective Search\\\normalfont\small $\sim$2k region proposals};

    % three candidate proposals of different size/position/color
    \draw[propedge, color=dlorange!90!black] ($(img2)+(-0.85,0.15)$) rectangle ++(0.95,0.75);
    \draw[propedge, color=dlred!85!black]    ($(img2)+(-0.15,-0.85)$) rectangle ++(1.05,0.65);
    \draw[propedge, color=dlgreen!55!black]  ($(img2)+(0.30,0.25)$) rectangle ++(0.65,0.60);

    \draw[flow] (img1) ++(\S/2,0) -- ++(0.9,0);

    % ===== Stage 3: warp each proposal to a fixed size =====
    \coordinate (warpanchor) at (8.0,0);

    \node[draw, propedge, color=dlorange!90!black, fill=dlorange!10, minimum width=0.85cm, minimum height=0.85cm] (w1) at ($(warpanchor)+(0,0.95)$) {};
    \node[draw, propedge, color=dlgreen!55!black,  fill=dlgreen!10,  minimum width=0.85cm, minimum height=0.85cm] (w2) at ($(warpanchor)+(0,0)$)    {};
    \node[draw, propedge, color=dlred!85!black,    fill=dlred!10,    minimum width=0.85cm, minimum height=0.85cm] (w3) at ($(warpanchor)+(0,-0.95)$) {};

    \node[lbl, align=center, below=0.55cm of w3, text width=2.3cm] {\bfseries\small Warp to fixed size\\\normalfont\small (e.g., $224\times224$)};

    \draw[flow, color=dlorange!90!black] ($(img2)+(0.10,0.525)$) -- (w1.west);
    \draw[flow, color=dlgreen!55!black]  ($(img2)+(0.95,0.55)$)  -- (w2.west);
    \draw[flow, color=dlred!85!black]    ($(img2)+(0.90,-0.525)$) -- (w3.west);

    % ===== Stage 4: CNN feature extraction (run independently per proposal) =====
    \node[box, minimum width=1.9cm, minimum height=2.6cm,
          fill=dlblue!16, draw=dlblue!70!black, text=dlblue!30!black,
          right=1.4cm of warpanchor, anchor=west] (cnn) {CNN};
    \node[lbl, below=0.7cm of cnn, font=\scriptsize, align=center, text=dlgray!70!black, text width=2cm] {shared weights\\run per proposal};

    \draw[flow] (w1.east) -- ++(0.5,0) |- ($(cnn.west)+(0,0.55)$);
    \draw[flow] (w2.east) -- (cnn.west);
    \draw[flow] (w3.east) -- ++(0.5,0) |- ($(cnn.west)+(0,-0.55)$);
    \node[lbl] at ($(cnn.north)+(0,0.5)$) {\bfseries Feature Extraction};

    % ===== Stage 5: per-class classification head =====
    \node[box, minimum width=2.6cm, minimum height=1.0cm,
          fill=dlgreen!14, draw=dlgreen!55!black, text=dlgreen!25!black,
          right=1.5cm of cnn.east, anchor=west] (svm) {SVM classifiers};

    \draw[flow] (cnn.east) -- (svm.west);

    \node[lbl, align=left, right=0.5cm of svm, font=\small] (out1)
        {person? \textcolor{dlgreen!45!black}{yes}\\aeroplane? \textcolor{dlred!60!black}{no}};

    \draw[flow] (svm.east) -- (out1.west);

\end{tikzpicture}
}
    \caption{The R-CNN pipeline: Selective Search extracts $\sim$2k region proposals from the input image, each of which is warped to a fixed size and passed independently through a CNN. The resulting features are then classified by per-class SVMs.}
    \label{fig:rcnn-architecture}
\end{figure}

This approach was a significant improvement over previous methods, but it still had a few major limitations, such as:
\begin{itemize}
    \item \textbf{NOT an end-to-end architecture}: The R-CNN architecture is not an end-to-end architecture, as it requires separate training for the region proposal, feature extraction, and classification steps.
    Deep learning models are known to perform better when trained end-to-end, as they can learn to optimize the entire process jointly.
    
    \item \textbf{Image resizing}: Each proposal must be resized to a fixed size (e.g., 224x224) before being fed into the CNN for feature extraction.
    This can lead to a loss of information, especially for small objects, and can also introduce distortions in the aspect ratio of the proposals.
    
    \item \textbf{Computationally expensive}: The R-CNN architecture is computationally expensive, as it requires running the CNN for each proposal, 
    which can be in the order of thousands per image (takes around 47 seconds per image on a GPU).
    
    \item \textbf{Selective Search}: The region proposal step relies on Selective Search, which is a hand-crafted and fixed algorithm.
    It may not be optimal for all types of images and objects, and can never be improved through learning.
\end{itemize}

\subsubsection{Fast R-CNN}
Fast R-CNN improves upon the original R-CNN by addressing its heavy computational redundancy and disjointed training pipeline. 
Instead of processing thousands of individual region proposals through a CNN, it utilizes a \textbf{shared computation strategy}.

The main innovations of Fast R-CNN are:
\begin{itemize}
    \item \textbf{Region Proposals (unchanged)}: Fast R-CNN still relies on Selective Search to generate the $\sim$2k region proposals for each image, exactly as in R-CNN.
    What changes is how these proposals are processed afterward.

    \item \textbf{Shared Convolutional Architecture}: Instead of running the CNN separately for each proposal,
    Fast R-CNN processes the entire image through a convolutional network to produce a global feature map. 
    
    \item \textbf{RoI Pooling Layer}: To handle the variable sizes of the region proposals, Fast R-CNN introduces a \textbf{Region of Interest} (RoI) \textbf{pooling} layer,
    which converts each region proposal into a fixed-size feature representation regardless of its original size. It works as follows:
    \begin{enumerate}
        \item Each region proposal is mapped onto the convolutional feature map computed for the entire image.
        \item The mapped region is divided into a fixed-size grid (e.g., $7 \times 7$).
        \item Max pooling is performed independently within each grid cell, rather than over the region as a whole.
        \item The pooled values from all cells are concatenated into a fixed-size descriptor (e.g., a $7 \times 7 \times C$ tensor, where $C$ is the number of feature channels), which is then fed to the fully connected layers.
    \end{enumerate}
    
    \item \textbf{Multi-Task Loss}: Fast R-CNN also introduces a multi-task loss function that combines classification and bounding box regression into a single optimization problem.
    This allows the model to learn both tasks simultaneously, improving the overall performance and efficiency of the training process.
\end{itemize}

\begin{figure}[H]
    \centering
    \resizebox{\textwidth}{!}{\usetikzlibrary{arrows.meta, calc, positioning}

\begin{tikzpicture}[
    >=Stealth,
    flow/.style={-{Stealth[length=7pt]}, line width=1pt},
    lbl/.style={font=\small, align=center},
    steplbl/.style={font=\small\bfseries, align=center, text=dlblue!35!black},
    box/.style={draw, thick, rounded corners, align=center, font=\bfseries\small},
]
    \def\photo#1#2#3{% #1=center coordinate name, #2=size, #3=fill
        \begin{scope}
            \clip (#1) ++(-#2/2,-#2/2) rectangle ++(#2,#2);
            \fill[#3] (#1) ++(-#2/2,-#2/2) rectangle ++(#2,#2);
            \fill[dlyellow!85!white] ($(#1)+(-0.30*#2,0.30*#2)$) circle (0.11*#2);
            \fill[dlgreen!45!black] ($(#1)+(-0.55*#2,-0.5*#2)$) -- ++(0.55*#2,0.42*#2) -- ++(0.55*#2,-0.42*#2) -- cycle;
            \fill[dlgreen!25!black] ($(#1)+(-0.05*#2,-0.5*#2)$) -- ++(0.38*#2,0.3*#2) -- ++(0.38*#2,-0.3*#2) -- cycle;
        \end{scope}
        \draw[draw=dlgray!80, thick] (#1) ++(-#2/2,-#2/2) rectangle ++(#2,#2);
    }

    \def\S{2.0}

    % ===== Stage 1: whole image through a single shared ConvNet =====
    \coordinate (img) at (0,0);
    \photo{img}{\S}{dlsky!12}
    \node[steplbl] at ($(img)+(0,\S/2+0.4)$) {Input Image};

    \node[box, minimum width=1.5cm, minimum height=2.5cm,
          fill=dlblue!16, draw=dlblue!70!black, text=dlblue!30!black,
          right=1.3cm of img] (convnet) {Deep\\ConvNet};
    \draw[flow] (img) ++(\S/2,0) -- (convnet.west);

    % ===== Stage 2: convolutional feature map for the whole image =====
    \coordinate (fmB) at ($(convnet.east)+(1.6,-0.9)$);
    \coordinate (fmA) at ($(convnet.east)+(2.5,0.9)$);
    \draw[tensorblock, line width=0.8pt] (fmB) rectangle (fmA);
    \foreach \x in {0,0.3,0.6,0.9} {\draw[tensoredge, thin] ($(fmB)+(\x,0)$) -- ($(fmB)+(\x,1.8)$);}
    \foreach \y in {0,0.3,...,1.8} {\draw[tensoredge, thin] ($(fmB)+(0,\y)$) -- ($(fmA)+(0,\y-1.8)$);}
    \draw[flow] (convnet.east) -- ($(fmB)!0.5!(fmA)$ |- convnet.east);
    \node[steplbl] at ($(fmA)!0.5!(fmB)+(0,1.05)$) {Feature Map};

    % ---- one RoI, projected from the original proposal onto the feature map ----
    \coordinate (roiFM) at ($(fmB)+(0.15,0.55)$); % bottom-left corner of the projected RoI
    \coordinate (roiFMright) at ($(roiFM)+(0.6,0.35)$);
    \coordinate (roiFMtopmid) at ($(roiFM)+(0.3,0.7)$);
    \draw[draw=dlorange!90!black, line width=1.1pt] (roiFM) rectangle ++(0.6,0.7);
    \draw[draw=dlorange!90!black, line width=1pt] ($(img)+(-0.25,-0.5)$) rectangle ++(0.9,0.95);

    % region proposals still come from Selective Search, unchanged from R-CNN
    \coordinate (imgbottom) at ($(img)+(0,-\S/2)$);
    \node[lbl, font=\scriptsize, align=center, text=dlorange!75!black,
          below=0.5cm of imgbottom, text width=3.4cm] (proplabel) {region proposal\\(Selective Search)};
    \draw[dlorange!80!black, thin] (proplabel.north) -- ($(img)+(0.2,-0.5)$);

    \coordinate (proj1) at ($(img)+(0.65,-0.025)$); % right edge of the proposal on the image
    \coordinate (proj2) at ($(proj1)+(0.5,0)$);      % step into the image/ConvNet gap
    \coordinate (proj3) at ($(proj2)+(0,1.70)$);     % rise above both the ConvNet and the feature map
    \coordinate (proj4) at (roiFMtopmid |- proj3);
    \draw[dlorange!80!black, dashed, line width=0.7pt] (proj1) -- (proj2) -- (proj3) -- (proj4) -- (roiFMtopmid);
    \node[font=\scriptsize, text=dlorange!70!black] at ($(proj3)!0.5!(proj4)+(0,0.22)$) {RoI projection};

    % ===== Stage 3: RoI pooling layer =====
    \node[box, minimum width=1.7cm, minimum height=1.6cm,
          fill=dlorange!14, draw=dlorange!75!black, text=dlorange!35!black,
          right=1.9cm of fmB, yshift=0.9cm, anchor=south west] (roipool) {RoI\\Pooling};
    \draw[flow, color=dlorange!85!black] (roiFMright) -- ++(0.3,0) |- (roipool.west);

    \node[draw, fill=dlorange!25, draw=dlorange!80!black, line width=0.8pt,
          minimum width=0.6cm, minimum height=0.6cm, right=1.0cm of roipool] (roivec) {};
    \draw[flow] (roipool.east) -- (roivec.west);
    \node[lbl, font=\scriptsize, below=0.1cm of roivec, align=center] {fixed-size\\RoI feature};

    % ===== Stage 4: fully connected layers =====
    \node[box, minimum width=1.1cm, minimum height=1.5cm,
          fill=dlblue!10, draw=dlblue!55!black, text=dlblue!30!black,
          right=0.9cm of roivec] (fc) {FC};
    \draw[flow] (roivec.east) -- (fc.west);

    % ===== Stage 5: two sibling output heads =====
    \node[box, minimum width=2.3cm, minimum height=0.9cm,
          fill=dlgreen!14, draw=dlgreen!55!black, text=dlgreen!25!black,
          right=1.3cm of fc.east, yshift=0.7cm, anchor=west] (cls) {softmax\\classifier};
    \node[box, minimum width=2.3cm, minimum height=0.9cm,
          fill=dlred!12, draw=dlred!70!black, text=dlred!45!black,
          right=1.3cm of fc.east, yshift=-0.7cm, anchor=west] (bbox) {bbox\\regressor};

    \draw[flow] (fc.east) -- ++(0.5,0) |- (cls.west);
    \draw[flow] (fc.east) -- ++(0.5,0) |- (bbox.west);

    \node[lbl, align=left, right=0.45cm of cls, font=\small] {class\\probabilities};
    \node[lbl, align=left, right=0.45cm of bbox, font=\small] {refined box\\coordinates};

    \node[lbl] at ($(roipool.north)!0.5!(fc.north)+(0,0.75)$) {\bfseries\small Multi-task loss};
    \draw[dlgray!50, dashed] ($(cls.north west)+(-0.35,0.35)$) rectangle ($(bbox.south east)+(0.35,-0.35)$);

\end{tikzpicture}
}
    \caption{The Fast R-CNN architecture: the whole image is passed once through a deep ConvNet to produce a shared feature map. Each region proposal is projected onto this feature map and reshaped by the RoI pooling layer into a fixed-size feature vector, which is then fed to fully connected layers ending in a softmax classifier and a bounding box regressor, jointly optimized with a multi-task loss.}
    \label{fig:fast-rcnn-architecture}
\end{figure}

This approach allowed reducing inference time significantly (to around 0.3 seconds per image on a GPU, from the 47 seconds of R-CNN), while also slightly improving the accuracy of the model.

\subsubsection{Faster R-CNN}
While Fast R-CNN significantly reduced the computational cost of the detection head, it remained bottlenecked by external region proposal 
methods like \textbf{Selective Search}, which typically took $\sim 2$ seconds per image on a CPU. 
Not only did this slow down the inference time, but it also prevented the model from being trained end-to-end, as the region proposal step was not learnable.

Faster R-CNN addresses this by introducing a \textbf{Region Proposal Network (RPN)}, allowing for a truly unified, end-to-end object detection framework.

The architecture of Faster R-CNN consists of the following components:
\begin{itemize}
    \item \textbf{Shared Convolutional Network}: Like its predecessor, it processes the entire image once to generate a high-level feature map.
    
    \item \textbf{RPN}: A small, fully convolutional network that slides over the feature map.
    At each sliding window location, it predicts objectness scores and bounding box offsets for $k$ \textbf{anchor boxes} of varying scales and aspect ratios.
    This implies that, for each position in the feature map, the RPN generates $k$ proposals, which are then filtered using non-maximum suppression (NMS) to remove redundant boxes.
    
    \item \textbf{RoI Pooling Layer}: This layer crops and reshapes the regions proposed by the RPN from the shared feature map into fixed-size feature vectors.
    
    \item \textbf{Classification and Bounding Box Regression}: The final stage consists of fully connected layers that predict the specific object class and perform fine-grained coordinates refinement.
\end{itemize}

\begin{figure}[H]
    \centering
    \resizebox{\textwidth}{!}{\usetikzlibrary{arrows.meta, calc, positioning}

\begin{tikzpicture}[
    >=Stealth,
    flow/.style={-{Stealth[length=7pt]}, line width=1pt},
    lbl/.style={font=\small, align=center},
    steplbl/.style={font=\small\bfseries, align=center, text=dlblue!35!black},
    box/.style={draw, thick, rounded corners, align=center, font=\bfseries\small},
]
    \def\photo#1#2#3{% #1=center coordinate name, #2=size, #3=fill
        \begin{scope}
            \clip (#1) ++(-#2/2,-#2/2) rectangle ++(#2,#2);
            \fill[#3] (#1) ++(-#2/2,-#2/2) rectangle ++(#2,#2);
            \fill[dlyellow!85!white] ($(#1)+(-0.30*#2,0.30*#2)$) circle (0.11*#2);
            \fill[dlgreen!45!black] ($(#1)+(-0.55*#2,-0.5*#2)$) -- ++(0.55*#2,0.42*#2) -- ++(0.55*#2,-0.42*#2) -- cycle;
            \fill[dlgreen!25!black] ($(#1)+(-0.05*#2,-0.5*#2)$) -- ++(0.38*#2,0.3*#2) -- ++(0.38*#2,-0.3*#2) -- cycle;
        \end{scope}
        \draw[draw=dlgray!80, thick] (#1) ++(-#2/2,-#2/2) rectangle ++(#2,#2);
    }

    \def\S{1.9}

    % ===== Stage 1: shared convolutional backbone =====
    \coordinate (img) at (0,0);
    \photo{img}{\S}{dlsky!12}
    \node[steplbl] at ($(img)+(0,\S/2+0.4)$) {Input Image};

    \node[box, minimum width=1.5cm, minimum height=2.3cm,
          fill=dlblue!16, draw=dlblue!70!black, text=dlblue!30!black,
          right=1.2cm of img] (convnet) {Shared\\ConvNet};
    \draw[flow] (img) ++(\S/2,0) -- (convnet.west);

    % ===== Stage 2: shared feature map =====
    \coordinate (fmB) at ($(convnet.east)+(1.2,-0.85)$);
    \coordinate (fmA) at ($(convnet.east)+(2.0,0.85)$);
    \draw[tensorblock, line width=0.8pt] (fmB) rectangle (fmA);
    \foreach \x in {0,0.4,0.8} {\draw[tensoredge, thin] ($(fmB)+(\x,0)$) -- ($(fmB)+(\x,1.7)$);}
    \foreach \y in {0,0.34,...,1.7} {\draw[tensoredge, thin] ($(fmB)+(0,\y)$) -- ($(fmA)+(0,\y-1.7)$);}
    \draw[flow] (convnet.east) -- ($(fmB)!0.5!(fmA)$ |- convnet.east);
    \node[steplbl] at ($(fmA)!0.5!(fmB)+(0,1.05)$) {Feature Map};

    % ===== Stage 3: Region Proposal Network branching upward =====
    \coordinate (fmtop) at ($(fmA)!0.5!(fmB)$);
    \node[box, minimum width=2.3cm, minimum height=1.0cm,
          fill=dlpurple!18, draw=dlpurple!65!black, text=dlpurple!45!black,
          above=1.55cm of fmtop] (rpn) {Region Proposal\\Network};
    \draw[flow, color=dlpurple!65!black] (fmtop) -- (rpn.south);

    % small anchor-box glyph next to the RPN, echoing the sliding-window/k-anchors idea
    \begin{scope}[xshift=0cm]
        \coordinate (anchoranchor) at ($(rpn.east)+(0.55,0)$);
        \draw[dlpurple!70!black, line width=0.7pt] ($(anchoranchor)+(-0.32,-0.22)$) rectangle ++(0.64,0.44);
        \draw[dlpurple!70!black, line width=0.7pt] ($(anchoranchor)+(-0.22,-0.32)$) rectangle ++(0.44,0.64);
        \draw[dlpurple!70!black, line width=0.7pt] ($(anchoranchor)+(-0.44,-0.14)$) rectangle ++(0.88,0.28);
        \node[lbl, font=\scriptsize, below=0.32cm of anchoranchor, align=center, text=dlpurple!45!black]
            {$k$ anchors\\per position};
    \end{scope}

    % ===== Stage 4: proposals produced by the RPN, projected back on the feature map =====
    \coordinate (propcenter) at ($(rpn.north)+(0,1.25)$);
    \draw[tensorblock alt, line width=0.8pt] ($(propcenter)+(-0.85,-0.7)$) rectangle ($(propcenter)+(0.85,0.7)$);
    \draw[draw=dlorange!90!black, line width=1pt] ($(propcenter)+(-0.55,-0.15)$) rectangle ++(0.6,0.55);
    \draw[draw=dlred!85!black, line width=1pt] ($(propcenter)+(-0.05,-0.45)$) rectangle ++(0.75,0.45);
    \draw[draw=dlgreen!55!black, line width=1pt] ($(propcenter)+(0.15,0.05)$) rectangle ++(0.5,0.45);
    \draw[flow, color=dlpurple!65!black] (rpn.north) -- (propcenter);
    \node[steplbl] at ($(propcenter)+(0,0.95)$) {Proposals};

    % ===== Stage 5: RoI pooling combines feature map + proposals =====
    \node[box, minimum width=1.7cm, minimum height=1.6cm,
          fill=dlorange!14, draw=dlorange!75!black, text=dlorange!35!black,
          right=3.7cm of fmB, yshift=0.85cm, anchor=south west] (roipool) {RoI\\Pooling};

    \draw[flow, color=dlorange!85!black] (propcenter) -- (propcenter -| roipool.north) -- (roipool.north);
    \draw[flow] (fmA)+(0,-0.2) -- ++(0.6,0) |- (roipool.west);

    \node[draw, fill=dlorange!25, draw=dlorange!80!black, line width=0.8pt,
          minimum width=0.6cm, minimum height=0.6cm, right=1.0cm of roipool] (roivec) {};
    \draw[flow] (roipool.east) -- (roivec.west);
    \node[lbl, font=\scriptsize, below=0.1cm of roivec, align=center] {fixed-size\\RoI feature};

    % ===== Stage 6: fully connected layers =====
    \node[box, minimum width=1.1cm, minimum height=1.5cm,
          fill=dlblue!10, draw=dlblue!55!black, text=dlblue!30!black,
          right=0.9cm of roivec] (fc) {FC};
    \draw[flow] (roivec.east) -- (fc.west);

    % ===== Stage 7: classification head =====
    \node[box, minimum width=2.3cm, minimum height=0.9cm,
          fill=dlgreen!14, draw=dlgreen!55!black, text=dlgreen!25!black,
          right=1.3cm of fc.east, anchor=west] (cls) {classifier};

    \draw[flow] (fc.east) -- (cls.west);

    \node[lbl, align=left, right=0.45cm of cls, font=\small] {object class};

\end{tikzpicture}
}
    \caption{The Faster R-CNN architecture: a shared ConvNet produces a feature map that feeds both the Region Proposal Network, which slides over it to output $k$ anchor-based proposals per position, and the RoI pooling layer, which pools the proposed regions into fixed-size features for the final classifier.}
    \label{fig:faster-rcnn-architecture}
\end{figure}

\paragraph{Performance}
By replacing Selective Search with the RPN, Faster R-CNN achieves an inference time of approximately \textbf{0.2 seconds} per image on a GPU. 
Furthermore, because the RPN is trained specifically for the task, it typically yields higher-quality proposals than traditional computer vision algorithms, leading to improved mAP.

\subsubsection{You Only Look Once}
\textbf{YOLO} (You Only Look Once) is a different approach to object detection. Instead of using a two-stage process like R-CNN and its variants,
which was very slow and hard to train, YOLO treats object detection as a single regression problem, 
directly predicting bounding boxes and class probabilities from the input image in one pass through the network.

The process of YOLO can be summarized as follows:
\begin{itemize}
    \item \textbf{Grid Division}: The input image is divided into an $S \times S$ grid. 
    Each grid cell is responsible for predicting $m$ bounding boxes;
    
    \item \textbf{Bounding Box Prediction}: For each bounding box, the model predicts the coordinates of
    the box and a class confidence score that indicates the likelihood of the box containing an object and the class of the object;
    
    \item \textbf{Bounding Box Selection}: During inference, the model selects only the bounding boxes which have
    a confidence score above a certain threshold.
\end{itemize}

This design allows YOLO to achieve real-time object detection, with inference times as low as 0.02 seconds per image on a GPU, while still maintaining competitive accuracy compared to two-stage detectors.
Even YOLO has some limitations, such as:

\begin{itemize}
    \item \textbf{Small Object Detection}: YOLO can struggle with detecting small objects, 
    as the grid cell responsible for predicting the object may not have enough information to accurately localize it.
    
    \item \textbf{Different Aspect Ratios}: YOLO uses a fixed set of anchor boxes, which may not be optimal for all object shapes and sizes, 
    leading to suboptimal performance for objects with unusual aspect ratios.
    
    \item \textbf{Loss Function}: The original YOLO loss is very complex and can be difficult to optimize,
    as it combines classification and localization losses, which can have different scales and gradients.
\end{itemize}

\subsection{Instance Segmentation}
A different task that can be tackled with CNNs is \textbf{instance segmentation}.
Instance segmentation is the task of assigning a class label to each pixel in an image, while also distinguishing between different instances of the same class.

For example, in an image containing multiple people, instance segmentation would not only identify the pixels belonging to the "person" class, 
but also differentiate between each individual person.

This task is more complex than object detection, as it requires not only localizing objects but also accurately segmenting them at the pixel level.

\subsubsection*{Mask R-CNN}
One of the most popular architectures for instance segmentation is \textbf{Mask R-CNN}, showing again the flexibility of CNNs to be adapted to different tasks.
Mask R-CNN builds upon the Faster R-CNN architecture by adding a branch for predicting segmentation masks on each Region of Interest (RoI).

To achieve this, Mask R-CNN introduces:
\begin{itemize}
    \item \textbf{RoIAlign Layer}: Instead of using the RoI Pooling layer from Faster R-CNN,
    which can quantize the RoI and lead to misalignments between the RoI and the extracted features,
    Mask R-CNN uses a RoIAlign layer that preserves the spatial alignment of the features, allowing for more accurate mask predictions.
    
    \item \textbf{Mask Branch}: In addition to the classification and bounding box regression branches, Mask R-CNN includes a separate branch that predicts a binary mask for each RoI.
    This branch is typically implemented as a small \textbf{fully convolutional network} (FCN) that takes the aligned features from the RoIAlign layer and outputs a mask of fixed size for each class, in a pixel-to-pixel manner that preserves the spatial layout of the RoI.
\end{itemize}

\begin{figure}[H]
    \centering
    \resizebox{\textwidth}{!}{\usetikzlibrary{arrows.meta, calc, positioning}

\begin{tikzpicture}[
    >=Stealth,
    flow/.style={-{Stealth[length=7pt]}, line width=1pt},
    lbl/.style={font=\small, align=center},
    steplbl/.style={font=\small\bfseries, align=center, text=dlblue!35!black},
    box/.style={draw, thick, rounded corners, align=center, font=\bfseries\small},
]
    \def\photo#1#2#3{% #1=center coordinate name, #2=size, #3=fill
        \begin{scope}
            \clip (#1) ++(-#2/2,-#2/2) rectangle ++(#2,#2);
            \fill[#3] (#1) ++(-#2/2,-#2/2) rectangle ++(#2,#2);
            \fill[dlyellow!85!white] ($(#1)+(-0.30*#2,0.30*#2)$) circle (0.11*#2);
            \fill[dlgreen!45!black] ($(#1)+(-0.55*#2,-0.5*#2)$) -- ++(0.55*#2,0.42*#2) -- ++(0.55*#2,-0.42*#2) -- cycle;
            \fill[dlgreen!25!black] ($(#1)+(-0.05*#2,-0.5*#2)$) -- ++(0.38*#2,0.3*#2) -- ++(0.38*#2,-0.3*#2) -- cycle;
        \end{scope}
        \draw[draw=dlgray!80, thick] (#1) ++(-#2/2,-#2/2) rectangle ++(#2,#2);
    }

    \def\S{1.9}

    % ===== Stage 1: shared convolutional backbone =====
    \coordinate (img) at (0,0);
    \photo{img}{\S}{dlsky!12}
    \node[steplbl] at ($(img)+(0,\S/2+0.4)$) {Input Image};

    \node[box, minimum width=1.5cm, minimum height=2.3cm,
          fill=dlblue!16, draw=dlblue!70!black, text=dlblue!30!black,
          right=1.2cm of img] (convnet) {Shared\\ConvNet};
    \draw[flow] (img) ++(\S/2,0) -- (convnet.west);

    % ===== Stage 2: shared feature map =====
    \coordinate (fmB) at ($(convnet.east)+(1.2,-0.85)$);
    \coordinate (fmA) at ($(convnet.east)+(2.0,0.85)$);
    \draw[tensorblock, line width=0.8pt] (fmB) rectangle (fmA);
    \foreach \x in {0,0.4,0.8} {\draw[tensoredge, thin] ($(fmB)+(\x,0)$) -- ($(fmB)+(\x,1.7)$);}
    \foreach \y in {0,0.34,...,1.7} {\draw[tensoredge, thin] ($(fmB)+(0,\y)$) -- ($(fmA)+(0,\y-1.7)$);}
    \draw[flow] (convnet.east) -- ($(fmB)!0.5!(fmA)$ |- convnet.east);
    \node[steplbl] at ($(fmA)!0.5!(fmB)+(0,1.05)$) {Feature Map};

    % ===== Stage 3: Region Proposal Network branching upward (inherited from Faster R-CNN) =====
    \coordinate (fmtop) at ($(fmA)!0.5!(fmB)$);
    \node[box, minimum width=2.3cm, minimum height=1.0cm,
          fill=dlpurple!18, draw=dlpurple!65!black, text=dlpurple!45!black,
          above=1.4cm of fmtop] (rpn) {Region Proposal\\Network};
    \draw[flow, color=dlpurple!65!black] (fmtop) -- (rpn.south);

    % ===== Stage 4: proposals produced by the RPN, projected back on the feature map =====
    \coordinate (propcenter) at ($(rpn.north)+(0,1.1)$);
    \draw[tensorblock alt, line width=0.8pt] ($(propcenter)+(-0.85,-0.6)$) rectangle ($(propcenter)+(0.85,0.6)$);
    \draw[draw=dlorange!90!black, line width=1pt] ($(propcenter)+(-0.55,-0.15)$) rectangle ++(0.6,0.5);
    \draw[draw=dlred!85!black, line width=1pt] ($(propcenter)+(-0.05,-0.4)$) rectangle ++(0.7,0.4);
    \draw[draw=dlgreen!55!black, line width=1pt] ($(propcenter)+(0.15,0.05)$) rectangle ++(0.5,0.4);
    \draw[flow, color=dlpurple!65!black] (rpn.north) -- (propcenter);
    \node[steplbl] at ($(propcenter)+(0,0.85)$) {Proposals};

    % ===== Stage 5: RoIAlign combines feature map + proposals =====
    \node[box, minimum width=1.7cm, minimum height=1.6cm,
          fill=dlorange!14, draw=dlorange!75!black, text=dlorange!35!black,
          right=3.7cm of fmB, yshift=0.85cm, anchor=south west] (roialign) {RoIAlign};
    \node[lbl, font=\scriptsize, below=0.15cm of roialign, align=center, text=dlgray!70!black, text width=3.2cm]
        {bilinear interpolation,\\no quantization};

    \draw[flow, color=dlorange!85!black] (propcenter) -- (propcenter -| roialign.north) -- (roialign.north);
    \draw[flow] (fmA)+(0,-0.2) -- ++(0.6,0) |- (roialign.west);

    \node[draw, fill=dlorange!25, draw=dlorange!80!black, line width=0.8pt,
          minimum width=0.6cm, minimum height=0.6cm, right=1.0cm of roialign] (roivec) {};
    \draw[flow] (roialign.east) -- (roivec.west);
    \node[lbl, font=\scriptsize, above=0.1cm of roivec, align=center] {aligned\\RoI feature};

    % ===== Stage 6: fully connected layers =====
    \node[box, minimum width=1.1cm, minimum height=1.5cm,
          fill=dlblue!10, draw=dlblue!55!black, text=dlblue!30!black,
          right=0.9cm of roivec] (fc) {FC};
    \draw[flow] (roivec.east) -- (fc.west);

    % ===== Stage 7: three sibling heads - classifier, bbox regressor, mask branch =====
    \node[box, minimum width=2.3cm, minimum height=0.85cm,
          fill=dlgreen!14, draw=dlgreen!55!black, text=dlgreen!25!black,
          right=1.3cm of fc.east, yshift=0.95cm, anchor=west] (cls) {classifier};
    \node[box, minimum width=2.3cm, minimum height=0.85cm,
          fill=dlred!12, draw=dlred!70!black, text=dlred!45!black,
          right=1.3cm of fc.east, anchor=west] (bbox) {bbox\\regressor};
    \node[box, minimum width=2.3cm, minimum height=0.85cm,
          fill=dlpurple!14, draw=dlpurple!65!black, text=dlpurple!45!black,
          right=1.3cm of fc.east, yshift=-0.95cm, anchor=west] (mask) {mask branch\\(small FCN)};

    \draw[flow] (fc.east) -- ++(0.5,0) |- (cls.west);
    \draw[flow] (fc.east) -- ++(0.5,0) -- (bbox.west);
    \draw[flow] (fc.east) -- ++(0.5,0) |- (mask.west);

    \node[lbl, align=left, right=0.45cm of cls, font=\small] {object class};
    \node[lbl, align=left, right=0.45cm of bbox, font=\small] {refined box\\coordinates};

    % small binary-mask glyph illustrating the mask branch's pixel-to-pixel output
    \coordinate (maskicon) at ($(mask.east)+(0.85,0)$);
    \draw[draw=dlgray!70, fill=dlgray!8, line width=0.7pt] ($(maskicon)+(-0.32,-0.32)$) rectangle ++(0.64,0.64);
    \fill[dlpurple!55!black] ($(maskicon)+(-0.14,-0.22)$) circle (0.14);
    \fill[dlpurple!55!black] ($(maskicon)+(-0.2,-0.3)$) rectangle ++(0.12,0.28);
    \draw[flow] (mask.east) -- ($(maskicon)+(-0.32,0)$);
    \node[lbl, align=left, right=0.3cm of maskicon, font=\small] {binary mask\\(per class)};

\end{tikzpicture}
}
    \caption{The Mask R-CNN architecture: it extends Faster R-CNN by replacing RoI Pooling with RoIAlign, which uses bilinear interpolation instead of quantization to preserve precise spatial alignment, and by adding a third sibling head, the mask branch, that predicts a binary mask for each RoI alongside the classifier and bounding box regressor.}
    \label{fig:mask-rcnn-architecture}
\end{figure}